\chapter{Getting started}
\label{ch:getting_started}
\index{Getting started}\index{First timetable}\index{Environment setup}

\epigraph{``The journey of a thousand miles begins with a single
step.''}{\textsc{Lao Tzu}, \emph{Tao Te Ching}, LXIV}

\lettrine[lines=3]{T}{his} chapter is a guided path that takes a
new reader --- whether a developer or an office clerk in a school
registry --- from a blank machine to the first timetable generated
and displayed on screen. The instructions are meant to be
copy-pasteable: if something goes wrong, the \emph{Common
pitfalls} box at the end of the chapter is the first place to
look.

\grokfig{fig_quickstart_5_passi}{The five steps of the quick
start, from installation to the first timetable in about thirty
minutes.}

\section{What we get}

By the end of the chapter you will have:
\begin{itemize}
\item a working local installation of $\pi$Tantum (FastAPI
  backend + SvelteKit frontend);
\item a test profile imported (\texttt{small}: 8 classes,
  19 teachers, 6 tracks);
\item a weekly timetable generated by the CP-SAT pipeline
  ``Phase A + Phase B per day'', visible in the \emph{Schedule}
  view;
\item a first custom constraint (written in the DSL) created from
  the Vincoli (Constraints) tab and re-applied to a new
  generation.
\end{itemize}

\sidenote{Estimated time for someone starting from scratch on
Windows 11: 30--40 minutes, of which 10 for downloading the
packages, 10 for the first CP-SAT generation, 10 for exploring
the UI.}

\section{Prerequisites}

You need three things installed and available in \texttt{\$PATH}:

\begin{itemize}
\item \textbf{Python 3.11 or later.} Check with
  \texttt{python --version}. On Windows the binary provided by the
  \emph{Microsoft Store}
  (\texttt{PythonSoftwareFoundation.Python.3.13}) works. On macOS
  it is best to install with \texttt{brew install python@3.13};
  on Linux the distribution package is usually enough.
\item \textbf{Node.js 20 or higher.} Check with
  \texttt{node --version}. It is needed for the SvelteKit
  dev-server of the UI. The exact installation instructions are
  in chapter~\ref{ch:installazione_en} (Installation).
\item \textbf{Git} to clone the repository and to be able to
  update the manual or the code without copying files by hand.
\end{itemize}

\section{Cloning and dependencies}

\sidenote{All the instructions assume a shell opened in the parent
directory where you want the repo to materialize. The paths in the
commands use forward slashes (\texttt{/}) for compatibility, but
on Windows they work just the same.}

\begin{lstlisting}[style=shell]
git clone https://github.com/gborghi/timetable.git
cd timetable

# (optional, recommended) dedicated virtual environment
python -m venv .venv
# activation: PowerShell
.\.venv\Scripts\Activate.ps1
# activation: bash/zsh
source .venv/bin/activate

# Python backend dependencies
pip install -r webui/backend/requirements.txt

# Node frontend dependencies
cd webui/frontend && npm install && cd ../..
\end{lstlisting}

The backend installation pulls in \texttt{ortools},
\texttt{fastapi}, \texttt{sqlalchemy}, \texttt{numpy}, \texttt{scipy}
and a handful of other packages. On the first run \texttt{pip}
may download $\sim$200 MB of binaries: this is normal.

\begin{erroricomunibox}
\textbf{``ortools wheel not found''}. On Python 3.13
\texttt{ortools} 9.15 is available for Windows, macOS and
Linux x86\_64. If the machine is ARM (Apple Silicon, Linux
ARM64) a more recent precompiled build or a source install may
be needed. The shortest way is
\texttt{pip install --upgrade ortools}.

\textbf{``Permission denied'' or a badly mounted virtualenv}.
On Windows, check that you launched PowerShell as a normal
user (not administrator) and that you have not set a
restrictive \texttt{ExecutionPolicy}: the command
\texttt{Get-ExecutionPolicy} must return at least
\texttt{RemoteSigned}.
\end{erroricomunibox}

\section{Starting the two processes}

\grokfig{shot_dashboard}{The Dashboard once the browser is open: at
the top, the guided first-steps checklist that ticks itself off as you
enter data.}

$\pi$Tantum runs as two cooperating processes:

\begin{itemize}
\item the FastAPI \textbf{backend} that serves the REST API and
  the CP-SAT logic;
\item the SvelteKit \textbf{frontend} that presents the user
  interface in the browser.
\end{itemize}

In two separate terminal tabs (or in two tmux panes):

\begin{lstlisting}[style=shell]
# Tab 1 (backend)
cd webui/backend
uvicorn main:app --reload --port 8000

# Tab 2 (frontend)
cd webui/frontend
npm run dev
\end{lstlisting}

The backend prints the message
\texttt{INFO: Uvicorn running on http://0.0.0.0:8000}; the
frontend in turn announces
\texttt{ready in $\sim$300ms} with the local URL.

\sidenote{Alternatively, the scripts
\texttt{webui/start.sh} (Linux/macOS),
\texttt{webui/start.bat} (Windows cmd) and
\texttt{webui/start.ps1} (PowerShell) start both processes in
one go. The \emph{Try it now} section below assumes this mode.}

\begin{esempiobox}[title=Try it now]
\begin{lstlisting}[style=shell]
cd webui
./start.ps1   # Windows PowerShell
# or
./start.sh    # Linux / macOS / Git Bash
\end{lstlisting}
If all goes well, after $\sim$10 seconds the default browser
opens on \texttt{http://localhost:5173}\,. The screen shows the
$\pi$Tantum dashboard with the navigation bar at the top and the
central \emph{Profili} (Profiles) panel.
\end{esempiobox}

\section{Importing the \texttt{small} profile}

The six test profiles (\texttt{small}, \texttt{medium},
\texttt{big}, \texttt{huge}, \texttt{superhuge}, \texttt{mega})
are pre-packaged as pickles in
\texttt{engine/scripts/data/<size>/}. The import copies them into
the SQLite database in use, recreating teachers, classes,
subjects, assignments and (if requested) the lessons of the
already-computed solution.

From the dashboard:

\begin{enumerate}
\item Click \emph{Importa profilo} (Import profile).
\item Select \texttt{small} from the drop-down menu.
\item Leave \emph{Importa lezioni dalla soluzione pre-calcolata}
  (Import lessons from the pre-computed solution) enabled if you
  want to see a timetable right away without re-running the
  solver. Disable it if you want to start from scratch.
\item Click \emph{Esegui import} (Run import).
\end{enumerate}

\begin{avvertenzabox}
The import \emph{wipes} the current database. Any data entered
manually is lost: save a dump beforehand if you need it. The
SQLite file is \texttt{webui/backend/timetable.db}.

The test pickles contain \emph{only the registry data}:
teachers, classes, tracks, assignments and (possibly) the
already-computed lessons. They do \emph{not} contain custom
constraints. After the import, the constraint tables
(unavailabilities, co-teaching, custom DSL constraints) will be
\textbf{empty}. If you want to reproduce a realistic scenario
with layered constraints, you must add them by hand from the
\emph{Docenti} (Teachers) tab (for the unavailabilities) and the
\emph{Vincoli} (Constraints) tab (for the logical and DSL
constraints).
\end{avvertenzabox}

\begin{biblioparvabox}
\sffamily
Under the hood, the endpoint
\texttt{POST /api/dataset/import-profile} calls
\texttt{optimization.import\_engine\_profile()} which in turn
delegates to three helpers in \texttt{engine\_io.py}:
\texttt{import\_school\_into\_db()} (registry data and classes),
\texttt{import\_profs\_into\_db()} (assignments) and ---
optionally ---
\texttt{import\_solution\_into\_db()} (lessons from
\texttt{solution\_<profile>\_optimized.pkl}).
\end{biblioparvabox}

\section{Viewing the timetable}

\grokfig{shot_orario}{The weekly timetable in the \emph{Per classe}
(by-class) view: each lesson is a draggable block; clicking one opens
its actions (move, delete, unschedule).}

Go to the \emph{Schedule} tab (\texttt{/schedule}). If you
imported the pre-computed solution you see a grid of
$36\text{ slots} = 6\text{ days} \times 6\text{ hours}$, with
one row per class, and in each cell the lesson: subject, teacher,
possibly a classroom. The colour coding (gold/ivory/indigo)
follows the vintage taste of the frontend.

The three essential commands:

\begin{itemize}
\item \textbf{Quick filters} (top): class, teacher, subject
  selection. When you filter by teacher you see only that
  teacher's slots; it is the handy view for whoever prepares a
  day's substitution.
\item \textbf{Drag-and-drop}: drag a lesson from one slot to
  another. If the move creates a hard violation the UI flashes
  red and cancels; if the solution stays valid the slot is
  updated instantly, and the \emph{Cost} panel at the top is
  refreshed.
\item \textbf{Export XLSX} (bottom right): generates an Excel file
  laid out according to the classic timetable form of Italian
  schools. Convenient to print or share by email.
\end{itemize}

\section{Running the pipeline from scratch}

\grokfig{shot_optimize}{The \emph{Ottimizza} (Optimize) page: the
\emph{Genera orario (consigliato)} (``Generate timetable,
recommended'') button on top for normal use; the raw solver knobs live
in the \emph{Modalit\`a avanzata} (advanced) accordion.}

Now let us do the round again, generating the timetable without
relying on the pre-computed solution.

\begin{enumerate}
\item Re-import \texttt{small} with \emph{Import lessons}
  disabled. The database is emptied of lessons; the Teacher,
  Class, Subject and Assignment tables stay populated.
\item Go to \emph{Optimize} (\texttt{/optimize}).
\item Choose the \emph{Phase A + Phase B (daily)} pipeline. Keep
  the default parameters: \texttt{time\_a=30s},
  \texttt{time\_b=12s}, workers $= 2$.
\item Click \emph{Avvia} (Start). The progress event appears at
  the bottom right: first the banner \emph{Phase A in progress\ldots}
  (during the weekly distribution of the daily counts), then six
  banners \emph{Phase B [day]} running one after the other.
\item When it finishes, the green banner \emph{Pipeline completed}
  automatically leads to the Schedule tab with the new timetable
  displayed.
\end{enumerate}

\sidenote{The distinction between \emph{Phase A} and \emph{Phase B}
is typical of $\pi$Tantum: the first distributes, for each
cattedra (a teacher's post: a teacher+class+subject triple with a
weekly hour count), the number of hours across the six days of the
week (without yet fixing the exact hours); the second, given the
per-day count, decides which time slot to assign to each lesson.
It is the same idea as the decomposition
\emph{schedule = day-distribution + day-assignment} historically
used in timetabling competitions.}

For larger profiles (\texttt{huge}, \texttt{mega}) the natural
pipeline is \emph{spectral decomposition} or
\emph{branch-and-price} with teacher granularity; chapter~\ref{ch:benchmarks_en}
offers the full matrix of times and chapter~\ref{ch:tecniche_avanzate_en}
explains when each combination makes sense.

\section{Adding a constraint from the UI}

\grokfig{shot_vincoli}{The \emph{Vincoli} (Constraints) page: add
rules (free days, unavailabilities, preferences) and run the
feasibility check.}

To close the loop we add a custom constraint and apply it. Open
the \emph{Vincoli} (Constraints) tab (\texttt{/constraints}) and
click \emph{+ Nuovo vincolo DSL} (+ New DSL constraint).

Example (literally pasteable):

\begin{lstlisting}
forall l in lessons where l.subject == Matematica:
    l.day != 6
\end{lstlisting}

Translated into words: ``for every mathematics lesson, let the
day not be Saturday''. Press \emph{Valida} (Validate) and then
\emph{Salva} (Save). The constraint enters as \texttt{level=hard}
in the \texttt{general\_constraints} table.

Go back to Optimize and re-run the pipeline. This time Phase A
notices the constraint through the hard DSL and redistributes the
mathematics hours over the first five days. If the constraint is
infeasible (for example: all the Saturdays were already the only
available space for mathematics) the system flags it and suggests
lowering it to \texttt{soft}.

\begin{esempiobox}[title=Try it now]
\begin{lstlisting}[style=shell]
# same constraint, launched via the API:
curl -X POST http://localhost:8000/api/constraints/general \
  -H 'Content-Type: application/json' \
  -d '{
    "name": "no_math_saturday",
    "expression": "forall l in lessons where l.subject == Matematica: l.day != 6",
    "level": "hard"
  }'
\end{lstlisting}
The JSON response contains the \texttt{id} of the created
constraint and \texttt{parsed\_ast\_json} with the validated AST.
To delete it: \texttt{DELETE /api/constraints/general/<id>}.
\end{esempiobox}

\section{Diagnostics and next steps}

At this point you can:

\begin{itemize}
\item explore the \emph{Diagnostica} (Diagnostics) tab
  (\texttt{/diagnostics}), which shows daily-load distributions
  per class and per teacher, sixth-hours, gaps, and the objective
  trend of Phase A;
\item read chapter~\ref{ch:dsl_generico_per_i_vincoli_en}
  (DSL: step-by-step tutorial) to write more sophisticated
  constraints;
\item take a look at chapter~\ref{ch:benchmarks_en} to understand
  which technique to choose if your school resembles one of the
  \texttt{huge} or \texttt{mega} profiles;
\item if you want to understand how $\pi$Tantum decides
  internally, open chapter~\ref{ch:metodo_cpsat_en} (CP-SAT) and the
  one on spectral decomposition.
\end{itemize}

\fleuron

\grokfig{fig_import_error_pesce}{The confused fish in front of a
file of incomplete format.}

\begin{erroricomunibox}
\textbf{``Backend unreachable''}. The dashboard tries to contact
\texttt{http://localhost:8000/api/health}; if it does not
respond, the red banner reports it. Check that the
\texttt{uvicorn} process is still running and that port 8000 is
not occupied by another service.

\textbf{``Pipeline stuck on Phase A''}. On \texttt{small}
Phase A finishes in 3--5 seconds; beyond 30 seconds suspect
that the database was imported badly (Assignment without
hours). Re-import the profile and try again.

\textbf{``DSL: parser error''}. The keywords are
\emph{case-sensitive}: \texttt{forall}, \texttt{where},
\texttt{in}, \texttt{exists}, \texttt{count} must be lowercase.
The identifiers (Borghi, Matematica, etc.) are literal values,
without quotes.

\textbf{``Port 5173 is already in use''}. It means that another
SvelteKit process has already started (perhaps from a previous
session that was not closed). On Linux/macOS:
\texttt{./webui/stop.sh} stops everything cleanly. On Windows:
close the cmd window left open or terminate the
\texttt{node.exe} process from the Task Manager.

\textbf{``The browser opens but I only see a blank page''}.
Open the browser console (\texttt{F12}): typically the problem
is that the backend is not responding. Check
\texttt{http://localhost:8000/api/health} in the browser: it
must return a small JSON \texttt{\{"status":"ok"\}}.

\textbf{``Pipeline completed but the timetable is empty''}.
Check on the Dashboard that the number of lessons is positive.
If it is zero, the pipeline probably failed silently: look at
the red messages in the log panel (Phase B).

\textbf{``I change the unavailability matrix but see no
effect''}. The changes take effect at the next pipeline run.
An edited matrix does not move lessons already placed: re-run
from \emph{Workflow}.
\end{erroricomunibox}
